% Figure: semi-infinite-solid temperature profiles in soil after a step to a
% sub-freezing surface, plotted at three elapsed times. The error-function
% profiles deepen with the square root of time; the depth where each curve
% crosses 0 degrees C is the frost penetration, advancing as sqrt(t). Used by
% the buried-pipe / frost-depth worked example. Coordinates are precomputed
% from T(x,t) = Ts + (T0 - Ts) erf( x / (2 sqrt(alpha t)) ) with Ts=-10,
% T0=8, alpha=5e-7 m^2/s at t = 10, 30, 90 days.
\begin{figure}[htbp]
\centering
\begin{tikzpicture}
\begin{axis}[
  width=12cm, height=7.5cm,
  xlabel={depth $x$ (m)},
  ylabel={soil temperature $T$ ($^{\circ}$C)},
  xmin=0, xmax=2.0, ymin=-12, ymax=10,
  axis lines=left,
  legend pos=south east,
  legend cell align=left,
  tick label style={font=\scriptsize},
  label style={font=\small},
  legend style={font=\scriptsize},
  clip=false,
]
\addplot[engblue, very thick] coordinates {
  (0.0,-10.0) (0.1,-8.458) (0.2,-6.935) (0.3,-5.446) (0.4,-4.007)
  (0.5,-2.634) (0.6,-1.338) (0.7,-0.128) (0.8,0.987) (0.9,2.004)
  (1.0,2.921) (1.1,3.737) (1.2,4.456) (1.3,5.082) (1.4,5.621)
  (1.5,6.079) (1.6,6.465) (1.7,6.785) (1.8,7.048) (1.9,7.262) (2.0,7.433)};
\addlegendentry{10 days}
\addplot[engamber, very thick] coordinates {
  (0.0,-10.0) (0.1,-9.109) (0.2,-8.221) (0.3,-7.339) (0.4,-6.468)
  (0.5,-5.611) (0.6,-4.769) (0.7,-3.947) (0.8,-3.147) (0.9,-2.371)
  (1.0,-1.622) (1.1,-0.901) (1.2,-0.21) (1.3,0.45) (1.4,1.078)
  (1.5,1.672) (1.6,2.234) (1.7,2.761) (1.8,3.255) (1.9,3.716) (2.0,4.145)};
\addlegendentry{30 days}
\addplot[engred, very thick] coordinates {
  (0.0,-10.0) (0.1,-9.485) (0.2,-8.971) (0.3,-8.458) (0.4,-7.947)
  (0.5,-7.438) (0.6,-6.933) (0.7,-6.432) (0.8,-5.935) (0.9,-5.443)
  (1.0,-4.957) (1.1,-4.477) (1.2,-4.004) (1.3,-3.538) (1.4,-3.08)
  (1.5,-2.63) (1.6,-2.189) (1.7,-1.757) (1.8,-1.334) (1.9,-0.92) (2.0,-0.517)};
\addlegendentry{90 days}
% zero-degree (frost) line
\draw[enggray, dashed] (axis cs:0,0) -- (axis cs:2.0,0);
\node[font=\scriptsize, enggray, anchor=south west] at (axis cs:1.45,0.2)
  {frost line ($0\,^{\circ}$C)};
\end{axis}
\end{tikzpicture}
\caption[Soil temperature profiles after a step to a sub-freezing
surface]{Soil temperature against depth after the surface is stepped to
$-10\,^{\circ}$C over ground initially at $+8\,^{\circ}$C, evaluated at ten,
thirty, and ninety days with the semi-infinite-solid error-function
solution and a soil diffusivity $\alpha = 5\times10^{-7}\,\text{m}^2/$s. The
profiles deepen with the square root of time, and the depth at which each
curve crosses $0\,^{\circ}$C is the frost penetration. The slow advance,
under a metre even after three months of hard cold, is why burying a water
line below the local frost line protects it and why the daily surface swing,
which has a far shorter penetration depth, never reaches that line at all.}
\label{fig:vol04ch05:semi-infinite-freeze}
\end{figure}
